\section{Prompt Templates (Expanded)}
\label{app:prompts}

The five task templates of Table~\ref{tab:tasks} are reproduced here
in their fully rendered form for unambiguous reproduction. The
\verb|<speech>| placeholder marks the position in the input
embedding stream where the merge step substitutes the projected
audio embedding sequence (Section~\ref{subsec:merge}); the
\verb|<start_speech>| and \verb|<end_speech>| tokens are emitted by
the chat template around every audio content item.

\begin{lstlisting}[style=plain, caption={Plain ASR prompt (no language hint, no hotwords).}]
Transcribe the following audio.<start_speech><speech><end_speech>
\end{lstlisting}

\begin{lstlisting}[style=plain, caption={Mandarin or English ASR with a language hint.}]
Transcribe the following audio.
Language: English<start_speech><speech><end_speech>
\end{lstlisting}

\begin{lstlisting}[style=plain, caption={Hotword-conditioned ASR; the language hint is optional.}]
Transcribe the following audio.
Language: Chinese
Hotwords: hw1, hw2, hw3<start_speech><speech><end_speech>
\end{lstlisting}

\begin{lstlisting}[style=plain, caption={Target-speaker ASR. The enrollment block precedes the mixed audio; the language and hotword lines are optional and follow the same convention as Hotword ASR.}]
Given the speaker's voice:<start_speech><speech><end_speech>
Transcribe what this speaker says in the following audio.<start_speech><speech><end_speech>
\end{lstlisting}

\section{Additional Results}
\label{app:more_results}

% TODO: Per-language WER, per-domain WER, more ablations that did not fit
% in the main text.

\section{Data Examples}
\label{app:data_examples}

% TODO: A handful of sanitised (audio identifier, transcript) examples per source.

\section{Retrieve Operator Constants}
\label{app:retrieve-constants}

Table~\ref{tab:retrieve_constants} lists the numerical constants of
the retrieval operator described in
Section~\ref{subsec:retrieve-algorithm}. The values are constrained
structurally by the channel-precision argument in
Equation~\eqref{eq:retrieve-score}: the threshold $\tau$ is chosen
so that any candidate supported only by low-precision channels has
a maximum attainable score below $\tau$ while a single
high-precision channel suffices to clear it. Specific values were
selected empirically and are reported here for reproducibility
rather than as ablated optima.

\begin{table}[ht]
  \caption{Retrieve operator constants for
    Section~\ref{subsec:retrieve-algorithm}.
    Channel notation: \texttt{sub} = exact substring;
    \texttt{py} = tone-stripped pinyin sliding window;
    \texttt{ovl} = character-set overlap;
    \texttt{jac} = bigram Jaccard;
    \texttt{wovl} = content-word overlap;
    \texttt{subsim} = best substring similarity under edit distance;
    \texttt{lcs} = longest-common-subsequence ratio.}
  \label{tab:retrieve_constants}
  \centering
  \small
  \begin{tabular}{llc}
    \toprule
    Route & Item & Value \\
    \midrule
    \multirow{6}{*}{CJK}
       & weight $w_{\texttt{sub}}$            & $2.0$ \\
       & weight $w_{\texttt{py}}$             & $1.5$ \\
       & weight $w_{\texttt{ovl}}$            & $0.5$ \\
       & weight $w_{\texttt{jac}}$            & $0.3$ \\
       & acceptance threshold $\tau$          & $1.0$ \\
       & CJK character ratio split            & $\ge 30\%$ characters CJK \\
    \midrule
    \multirow{8}{*}{Alphabetic}
       & weight $w_{\texttt{sub}}$            & $3.0$ \\
       & weight $w_{\texttt{wovl}}$           & $2.0$ \\
       & weight $w_{\texttt{subsim}}$         & $1.5$ \\
       & weight $w_{\texttt{lcs}}$            & $0.5$ \\
       & short-word penalty $\mathrm{pen}(|h|\!\le\!3 / \le\!4 / \le\!5)$ & $0.35 / 0.20 / 0.10$ \\
       & default subsim / lcs thresholds       & $0.65 / 0.60$ \\
       & length-5 fallback thresholds          & $0.75 / 0.70$ \\
       & multi-word strict thresholds          & $\mathrm{subsim}\ge 0.80,\ \mathrm{lcs}\ge 0.75$ \\
    \midrule
    Either & discriminative-gap factor       & $\mathrm{top1}/\mathrm{top10} \ge 1.30$ \\
    \midrule
    \multirow{2}{*}{Training augmentation}
       & ground-truth-miss probability $p_\mathrm{miss}$ & $0.12$ \\
       & hotword-line drop probability $p_\mathrm{drop}$ & $0.20$ \\
    \bottomrule
  \end{tabular}
\end{table}

\section{GRPO Hyperparameters}
\label{app:grpo-hparams}

Table~\ref{tab:grpo_hparams} lists the hyperparameters of the GRPO
recipe of Section~\ref{subsec:grpo}.

\begin{table}[ht]
  \caption{GRPO hyperparameters for
    Section~\ref{subsec:grpo}. Equation references point to the
    relevant terms in Section~\ref{subsec:grpo}.}
  \label{tab:grpo_hparams}
  \centering
  \small
  \begin{tabular}{lc}
    \toprule
    Item & Value \\
    \midrule
    Group size $G$ (Eq.~\ref{eq:grpo-advantage})       & $16$ \\
    Sampling temperature                                & $0.6$ \\
    PPO clip $\varepsilon$ (Eq.~\ref{eq:grpo-objective}) & $0.2$ \\
    KL coefficient $\beta$ (Eq.~\ref{eq:grpo-objective}) & $1\!\times\!10^{-3}$ \\
    Peak learning rate                                  & $1\!\times\!10^{-6}$ \\
    Epochs                                              & $1$ \\
    Per-device prompt batch                             & $8$ \\
    Gradient accumulation                               & $2$ \\
    ASR-accuracy reward weight $w_\mathrm{asr}$ (Eq.~\ref{eq:reward-total})   & $1.0$ \\
    Hotword-match reward weight $w_\mathrm{hw}$ (Eq.~\ref{eq:reward-total})   & $0.3$ \\
    \bottomrule
  \end{tabular}
\end{table}
